(1) 质量 \(m\) 的无自旋带电粒子电荷 \(q\)，被限制在二维半径 \(R\) 的无穷高圆盘势阱 \(r<R\)，边界为 Dirichlet 条件。无磁场时从极坐标 Schrödinger 方程严格求全部定态的本征能量、归一化波函数及角向简并性，记 \(j_{\ell,n}\) 为 \(J_\ell\) 的第 \(n\) 个正零点。
(2) 对角量子数为整数 \(\ell\) 的定态，求粒子数概率流密度及总 \(z\) 角动量；说明 \(\ell\) 和 \(-\ell\) 的态在无场下为何等能，并核验 \(\ell=0\) 的概率流。
(3) 现在施加垂直于圆盘的弱均匀磁场 \(B\)，取对称规范 \(\mathbf A=(Br/2)\mathbf e_\varphi\)。从最小耦合哈密顿量精确展开到 \(B^2\)，计算 \(\ell\ne0\) 的一阶 Zeeman 型分裂并讨论电子 \(q=-e\) 时的正负号。
(4) 对任意给定 \((\ell,n)\) 推导径向期望值 \(\langle r^2\rangle\) 的闭式 Bessel 零点公式，给出由 \(B^2\) 项导致的首个抗磁能量修正，并对基态 \(\ell=0\) 验证其为正值且 \(B\to0\) 时恢复零场谱；明确微扰阶数及忽略的自旋效应。
